\section{Introduction}\label{introduction}

\section{Introduction}\label{introduction-1}

Tabular data remain one of the most common forms of structured data used in practical machine learning applications. Financial records, demographic information, customer interactions, medical measurements, and many other real-world datasets are naturally represented as rows of observations and columns of heterogeneous features. Despite the rapid development of deep learning and foundation models in areas such as computer vision and natural language processing, conventional machine-learning methods remain highly competitive for tabular prediction tasks {[}1,2{]}.

In particular, gradient-boosted decision tree (GBDT) methods have established themselves as strong general-purpose approaches for tabular classification {[}2,3{]}. XGBoost is a scalable tree-boosting system designed for efficient and accurate learning {[}4{]}, while LightGBM introduced techniques such as Gradient-based One-Side Sampling and Exclusive Feature Bundling to improve the efficiency of gradient-boosted decision trees {[}5{]}. CatBoost introduced ordered boosting and specialized processing for categorical features to address prediction shift associated with conventional boosting approaches {[}6{]}.

These methods have remained particularly competitive in tabular learning. Comparative studies have shown that tree-based methods can outperform a range of deep learning approaches on medium-sized tabular datasets, while also providing practical advantages in training and optimization efficiency {[}2,3{]}. This continuing strength of gradient-boosted tree methods provides an important baseline when evaluating newer neural and foundation-model approaches for tabular data.

Although these methods remain highly competitive, their conventional learning paradigm requires a model to learn a new predictive function from the available training data for each task. This creates a particular challenge when only a limited number of labelled observations are available.

Recent developments in tabular foundation models have introduced a different approach. Instead of learning a new model entirely from scratch for each dataset, a foundation model can acquire general predictive capabilities during large-scale pre-training and subsequently apply those capabilities to previously unseen tasks. Tabular Prior-data Fitted Networks (TabPFN) represent one such approach. TabPFN uses a transformer-based architecture trained across large collections of synthetic tabular prediction tasks and applies the resulting learned algorithm to new tabular problems through in-context learning {[}7{]}.

The original TabPFN study reported particularly strong performance on small- and medium-sized tabular datasets, including datasets with up to approximately 10,000 samples, and compared TabPFN against established baselines including gradient-boosted tree models {[}7{]}. These findings raise an important practical question: whether the relative advantage of a tabular foundation model persists across different amounts of available training data and across different types of tabular datasets.

This question is particularly relevant because the practical choice between a foundation model and conventional tree-based methods should not be determined solely by performance on a single dataset or at a single training size. A model that performs particularly well when labelled data are scarce may have a different relative advantage when substantially more training observations become available. Similarly, performance differences may depend on the characteristics of the dataset and the evaluation metric being considered {[}2,3{]}.

The present study therefore conducts an empirical comparison of TabPFN with three established gradient-boosted tree algorithms: XGBoost, LightGBM, and CatBoost. The comparison is performed across three tabular binary-classification datasets and a range of training-data sizes. Multiple random seeds are used at each training size to reduce the influence of sampling variation. In addition to predictive performance, the study considers computational measurements in order to examine the practical trade-offs associated with the different approaches.

The primary research question is:

\begin{quote}
How does TabPFN compare with established gradient-boosted tree algorithms across different training-data regimes in tabular classification?
\end{quote}

This question is addressed through four supporting research questions:

\begin{enumerate}
\def\labelenumi{\arabic{enumi}.}
\tightlist
\item
  Does TabPFN provide an advantage when training data are limited?
\item
  How does the relative performance of TabPFN change as training size increases?
\item
  Is the performance advantage of TabPFN consistent across datasets and evaluation metrics?
\item
  What computational trade-offs exist between TabPFN and gradient-boosted tree models?
\end{enumerate}

Based on these questions, the study investigates four corresponding hypotheses. First, TabPFN is expected to be competitive or superior to gradient-boosted tree models when training data are limited. Second, its relative advantage is expected to change as more training data become available. Third, the magnitude of the difference between TabPFN and tree-based models is expected to vary across datasets and evaluation metrics. Finally, differences in predictive performance are expected to be accompanied by differences in computational cost.

The main contribution of this study is an empirical, training-data-aware comparison between a tabular foundation model and established gradient-boosted tree algorithms. Rather than evaluating the models only at a single dataset size, the study examines their behaviour across multiple training-data regimes and repeated random seeds. This allows the analysis to distinguish between overall model performance and performance under limited-data conditions.

A second contribution is the inclusion of multiple complementary evaluation dimensions. Predictive performance is assessed using Accuracy, Precision, Recall, F1-score, and ROC-AUC, while training and prediction time are used to characterize computational behaviour. This provides a broader basis for model comparison than relying on a single predictive metric.

Finally, the study aims to provide a balanced assessment of tabular foundation models. Rather than assuming that a newer foundation-model approach should replace established tree-based algorithms, the analysis investigates the conditions under which TabPFN provides an advantage and the conditions under which conventional gradient-boosted trees remain competitive. The resulting comparison is intended to provide practical evidence for selecting models according to dataset characteristics and available training data.
